\section{Qwen3.5 APPS replication protocol}
\label{app:qwen35-apps}

This replication keeps the APPS task fixed while changing the model. It uses
the same 1,444 honest--backdoor pairs as the Qwen2.5 and Gemma studies, split by
problem ID into 500 fit, 200 calibration, and 744 locked-test pairs. Both
members of a pair remain in the same split. The prompts, score tokens,
controller, audit budgets, resampling population, and realized-violation loss
also remain fixed. Thinking is disabled for every Qwen3.5 view.

The design was chosen after the Qwen2.5 and Gemma results were known. Before
reading any Qwen3.5 monitor outcome, we pinned the model and official SAE
revisions, declared the primary contrast, and sealed the source manifest. The
primary contrast is neutral-residual minus verdict-adjacent-residual mean
violations at 1\% attack prevalence and a 2\% audit budget. This is a local
prospective seal, not an independent preregistration.

The pinned model revision is
\path{c202236235762e1c871ad0ccb60c8ee5ba337b9a}; the pinned SAE revision is
\path{7bb2370d360e566b2d8acb8b3d15a0b2b88f52a8}. The source files and sealed
manifest ship with the artifact, so the aggregate scorecard can be tied to the
exact code and configuration used for the run without publishing per-example
evaluator labels.

We measure post-layer residuals at layers 12, 20, and 28. Fit data train each
probe; calibration chooses the layer and ridge value; the locked test is read
only after selection. One fit pair exceeds the 4,096-token limit and is
excluded under the frozen rule. No calibration or locked-test pair reaches the
limit. A compatibility-only smoke first checks prompt lengths, score tokens,
residual shapes, and the three SAE files without fitting or ranking an
observer.

The official-feature control uses 65,536-wide TopK-50 Qwen-Scope SAEs at the
same three layers. Calibration selects layer 12 and ridge 100. The SAEs were
trained on Qwen3.5-9B-Base, whereas the dense residuals and outcomes come from
post-trained Qwen3.5-9B. We declared this mismatch before the run. It makes the
SAE row a test of a released feature family under transfer, not a
checkpoint-matched comparison and not evidence that SAE features generally
fail for safety monitoring.
